\section{Thermodynamic and Column Formulation}\label{sec:model}
\subsection{System definition and modeling assumptions}
% P:model-framework-01
The formulation describes steady countercurrent absorption in a single packed bed of height $H$ and cross-sectional area $A$.
The axial coordinate $z$ increases from the gas inlet at $z=0$ to the liquid inlet at $z=H$.
Phase flow rates are positive in their respective directions of motion.
The gas contains CO$_2$, H$_2$O, N$_2$, and O$_2$; the apparent liquid components are CO$_2$, MEA, and H$_2$O.
CO$_2$ and water transfer between phases, while MEA, N$_2$, and O$_2$ are treated as nontransferring components.
The one-bed formulation neglects axial dispersion, heat loss to the surroundings, and pressure variation over the packed height.

% P:model-framework-02
The ePC-SAFT bulk liquid and equilibrium-manifold film share the nine-species description defined below.
Thermodynamic comparators use their specified species and reaction systems, with feeds expressed on the same analytical CO$_2$, MEA, and water basis.
The common absorber balances compare capture, loading, fugacity, and temperature without requiring identical internal species lists.
Local reaction equilibrium applies within the liquid, whereas the bulk gas and liquid remain separated by finite transfer resistances.
The CO$_2$ film is isothermal at the local liquid temperature, and its mobility constraints separate CO$_2$ transport from the water-transfer calculation.
This approximation combines local reaction equilibrium with constrained diffusive mobility; it does not resolve finite reaction rates or a separate temperature field inside a general multicomponent film.
\Cref{fig:absorber-calculation-sequence} summarizes the relationships between thermodynamics, transfer, and the column balances.

\begin{figure}[pos=htbp]
\centering
% Scientific dependency diagram; no implementation iteration or result status.
\begin{tikzpicture}[node distance=9mm and 14mm,
 cse stage/.append style={text width=49mm,minimum width=0pt},
 cse detail/.append style={text width=49mm},
 cse solver/.append style={text width=49mm,minimum width=0pt}]
\node[cse stage] (state) {Local column state\\Component flows, energy, pressure};
\node[cse stage,below=of state] (thermo) {Reactive thermodynamics\\Species, density, enthalpy, fugacity};
\node[cse stage,below=of thermo] (transfer) {Interfacial transfer\\CO$_2$, water, and energy fluxes};
\node[cse stage,below=of transfer] (balances) {Countercurrent column balances\\Gas and liquid inlet conditions};
\node[cse solver,below=of balances] (solve) {Boundary-value solution\\Conservative formulation and comparison methods};
\node[cse detail,right=of thermo] (path) {Equilibrium path\\Chemical-potential derivatives and constrained mobility};
\node[cse detail,below=of path] (film) {Local CO$_2$ film\\Path quadrature and scalar interface balance};
\foreach \a/\b in {state/thermo,thermo/transfer,transfer/balances,balances/solve}
 \draw[cse arrow] (\a) -- (\b);
\draw[cse link] (thermo.east) -- (path.west);
\draw[cse link] (path) -- (film);
\draw[cse link] (film.west) -- (transfer.east);
\draw[cse arrow] (solve.west) -- ++(-8mm,0) |- (state.west);
\end{tikzpicture}

\caption{Thermodynamic and transport structure of the absorber formulation. A common reactive liquid state supplies composition, thermodynamic properties, and the equilibrium path used in the local film calculation. Interfacial material and energy transfer couple the countercurrent phase balances.}
\label{fig:absorber-calculation-sequence}
\end{figure}

\subsection{Species, conserved components, and reactive equilibrium}\label{subsec:equilibrium}
% P:model-framework-03
The ePC-SAFT liquid species are CO$_2$, MEA, H$_2$O, MEAH$^+$, MEACOO$^-$, HCO$_3^-$, CO$_3^{2-}$, H$_3$O$^+$, and OH$^-$.
The independent reactions are
\begin{align}
2\ce{H2O}&\rightleftharpoons\ce{H3O+}+\ce{OH-},\\
\ce{CO2}+2\ce{H2O}&\rightleftharpoons\ce{HCO3-}+\ce{H3O+},\\
\ce{HCO3-}+\ce{H2O}&\rightleftharpoons\ce{CO3^2-}+\ce{H3O+},\\
\ce{MEACOO-}+\ce{H2O}&\rightleftharpoons\ce{MEA}+\ce{HCO3-},\\
\ce{MEAH+}+\ce{H2O}&\rightleftharpoons\ce{MEA}+\ce{H3O+}.
\label{eq:reactive-reactions}
\end{align}
For species amounts $\mathbf n$ and stoichiometric matrix $\boldsymbol\nu$, the conserved totals satisfy
\begin{align}
\mathbf B\mathbf n&=\mathbf b,&\mathbf B\boldsymbol\nu&=\mathbf0,
\label{eq:conserved-totals}\\
\sum_i\nu_{ir}\ln a_i&=\ln K_r(T).
\label{eq:reactive-equilibrium}
\end{align}
The rows of $\mathbf B$ specify the independent material totals imposed by the equilibrium formulation, with $\mathbf b=\mathbf B\mathbf n_{\mathrm{feed}}$ on the chosen feed-amount basis.
An amount vector in moles gives conserved totals in moles of the corresponding material quantities; the same transformation applies to molar flow vectors.
Electroneutrality is $\mathbf z^{\mathsf T}\mathbf n=0$, where $\mathbf z$ is the species charge vector.
% CHECKLIST-BEGIN: method-01
% P:model-framework-04
In the species order above, the selected two-row basis is
\begin{align}
\mathbf B=\begin{pmatrix}
1&2&0&2&3&1&1&0&0\\
0&1&0&1&1&0&0&0&0
\end{pmatrix}.
\label{eq:conserved-basis-matrix}
\end{align}
The rows count carbon atoms and amine nitrogen atoms, respectively; $\mathbf B$ is dimensionless and has size $2\times9$.
For apparent feed amounts $(n_C^{\mathrm{app}},n_A^{\mathrm{app}},n_W^{\mathrm{app}})$, define $n_0=n_C^{\mathrm{app}}+n_A^{\mathrm{app}}+n_W^{\mathrm{app}}$ and $\widehat n_j=n_j^{\mathrm{app}}/n_0$.
The normalized feed is $(\widehat n_C,\widehat n_A,\widehat n_W,0,0,0,0,0,0)^{\mathsf T}$ and its conserved totals are $(\widehat n_C+2\widehat n_A,\widehat n_A)^{\mathsf T}$ on the unit apparent-feed amount basis.
Multiplication by $n_0$ restores extensive amounts; applying the same construction to feed rates gives conserved flows in \unit{mol\,s^{-1}}.
The absorbed-CO$_2$ component total is $b_1-2b_2$, and the apparent loading is $(b_1-2b_2)/b_2$.
The complete equilibrium constraints also include total mass and electroneutrality:
\begin{align}
\mathbf M^{\mathsf T}\mathbf n&=m_0,
&m_0&=\mathbf M^{\mathsf T}\mathbf n_{\mathrm{feed}},\\
\mathbf z^{\mathsf T}\mathbf n&=0,
&\boldsymbol\beta&=(m_0,b_1,b_2)^{\mathsf T}.
\label{eq:complete-invariants}
\end{align}
Here $M_i$ is the species molar mass in \unit{kg\,mol^{-1}}; $m_0$ is in kilograms on the selected amount basis, and $\mathbf z$ is the charge vector in \cref{eq:stoichiometric-matrix}.
The mass row is additional to the two supplied material rows, and the charge total remains fixed at zero.
Water therefore affects both the conserved mass and the solvent dilution, even though it has no separate row in $\mathbf B$.
The mass and material totals are varied together when the apparent feed changes.
% CHECKLIST-END: method-01
Activities and equilibrium constants use a common standard-state convention for each reaction.
The aqueous solute convention uses a standard molality of \qty{1}{mol\,kg^{-1}} and water as reference solvent.

% P:model-framework-05
Apparent component amounts specify the feed to the equilibrium calculation; reaction products have zero amount in that physical feed representation.
The equilibrium species distribution is determined by \cref{eq:reactive-equilibrium} at the prescribed temperature, pressure, and conserved totals.
True mole fractions and concentrations are
\begin{align}
x_i=\frac{n_i}{\sum_j n_j},\qquad c_i=\rho_l x_i,
\label{eq:true-concentration}
\end{align}
where $\rho_l$ is the true-species molar density.
The apparent loading is $\alpha=F_C^l/F_A^l$, with $C$ denoting the CO$_2$ component and $A$ the MEA component.
Apparent component fractions describe feed and column balances; true-species fractions describe equilibrium and molecular transport.

\subsubsection{Reaction stoichiometry and standard states}
% P:model-framework-06
In the species order stated above and reaction order of \cref{eq:reactive-reactions}, the stoichiometric matrix is
\begin{align}
\boldsymbol\nu=
\begin{pmatrix}
0&-1&0&0&0\\
0&0&0&1&1\\
-2&-2&-1&-1&-1\\
0&0&0&0&-1\\
0&0&0&-1&0\\
0&1&-1&1&0\\
0&0&1&0&0\\
1&1&1&0&1\\
1&0&0&0&0
\end{pmatrix},\qquad
\mathbf z=(0,0,0,1,-1,-1,-2,1,-1)^{\mathsf T}.
\label{eq:stoichiometric-matrix}
\end{align}
Each reaction conserves charge, so $\mathbf z^{\mathsf T}\boldsymbol\nu=\mathbf0$.
The species amounts may change substantially with loading even though the conserved component totals remain fixed within each equilibrium calculation.
Reaction equilibrium is expressed through dimensionless activities; the numerical values of $K_r$ are consequently tied to the molality or solvent convention used to define those activities.

% P:model-framework-07
For a solute on the aqueous molality convention,
\begin{align}
a_i=\gamma_i^{(m)}\frac{m_i}{m^\circ},\qquad
\mu_i=\mu_i^\circ+RT\ln a_i,
\label{eq:molality-activity}
\end{align}
with $m^\circ=\qty{1}{mol\,kg^{-1}}$.
The solute molality is $m_i=n_i/(M_W n_W)=c_i/(M_Wc_W)$ in \unit{mol\,kg^{-1}}, where $M_W$ is the molar mass of water in \unit{kg\,mol^{-1}} and $n_W$ is the molecular-water amount in the equilibrium state.
Water has its solvent reference state and is excluded from this solute-molality definition.
The reaction standard chemical potential satisfies
\begin{align}
\sum_i\nu_{ir}\mu_i^\circ=-RT\ln K_r.
\end{align}
This relation connects the equilibrium constants to the chemical-potential convention used in the film.
A change in standard state requires the corresponding change in the equilibrium constant, rather than a change in the physical equilibrium.

% P:model-framework-08
The retained temperature correlations for the selected reaction system take the forms
\begin{align}
\ln K_r(T)=
\begin{cases}
a_r+b_r/T+c_r\ln T+s_r,&r=1,2,3,\\
a_4+b_4/T,&r=4,\\
-\ln(10)(A_5/T+B_5+C_5T),&r=5,
\end{cases}
\label{eq:reaction-constants}
\end{align}
where $T$ is evaluated numerically in kelvin and the coefficient units follow that convention.
The offsets $s_r$ express the standard-state conversion for the corresponding source relations.
Their derivatives retain the same temperature dependence:
\begin{align}
\odv{\ln K_r}{T}&=-\frac{b_r}{T^2}+\frac{c_r}{T},\qquad r=1,2,3,\\
\odv{\ln K_4}{T}&=-\frac{b_4}{T^2},\\
\odv{\ln K_5}{T}&=\ln(10)\left(\frac{A_5}{T^2}-C_5\right).
\end{align}
These derivatives enter the temperature response of the equilibrium species distribution.
A source correlation with an additional $d_rT$ term contributes $d_r$ to its first derivative and retains that term in both equilibrium and caloric evaluations.

\subsection{ePC-SAFT and comparison thermodynamics}\label{subsec:thermodynamics}
% P:model-framework-09
For ePC-SAFT, the dimensionless residual Helmholtz energy is
\begin{align}
a^{\mathrm{res}}&=a^{\mathrm{hc}}+a^{\mathrm{disp}}+a^{\mathrm{assoc}}+a^{\mathrm{DH}}+a^{\mathrm{Born}},
\label{eq:helmholtz}\\
P&=\rho RT\left[1+\rho\left(\pdv{a^{\mathrm{res}}}{\rho}\right)_{T,\mathbf x}\right],
\label{eq:eos-pressure}\\
\ln\phi_i&=\left(\pdv{(na^{\mathrm{res}})}{n_i}\right)_{T,V,n_{j\ne i}}-\ln Z,
\qquad Z=\frac{P}{\rho RT},\\
f_i&=x_i\phi_i P.
\label{eq:eos-fugacity}
\end{align}
Here $a^{\mathrm{res}}=A^{\mathrm{res}}/(nRT)$ and $\rho$ is molar density.
The terms represent hard-chain, dispersion, association, Debye--H\"uckel, and Born contributions \cite{Gross2001,Cameretti2005,Zuber2014}.
The equation of state supplies phase density, fugacity, chemical potentials, and the derivatives used to follow the equilibrium manifold \cite{Polley2026epcsaft}.
Gas properties use the complete CO$_2$/H$_2$O/N$_2$/O$_2$ mixture.

% P:model-framework-10
The Henry-law reference expresses dissolved molecular CO$_2$ fugacity as
\begin{align}
f_C^l=H_C^c(T,\mathbf x)c_C,
\label{eq:henry}
\end{align}
where $H_C^c$ has units of \unit{Pa\,m^3\,mol^{-1}} and $c_C$ is molecular CO$_2$ concentration.
The eNRTL comparison obtains activities from an excess-Gibbs-energy description \cite{Zhang2011}.
Each comparator uses internally consistent reaction constants and standard-state transformations.
Differences between complete thermodynamic models include any differences in their species and reaction descriptions, as specified in \cref{sec:methods}.
The common-property and full-property comparisons are specified separately in \cref{sec:methods}.

% P:model-enrtl-species
The eNRTL comparator contains six liquid species: CO$_2$, MEA, H$_2$O, MEAH$^+$, MEACOO$^-$, and HCO$_3^-$.
Its two combined formation reactions are
\begin{align}
\ce{CO2}+2\ce{MEA}&\rightleftharpoons\ce{MEAH+}+\ce{MEACOO-},\nonumber\\
\ce{CO2}+\ce{MEA}+\ce{H2O}&\rightleftharpoons\ce{MEAH+}+\ce{HCO3-}.
\label{eq:enrtl-combined-reactions}
\end{align}
The six-species comparator omits the explicit carbonate, hydronium, and hydroxide species in the nine-species ePC-SAFT formulation.
The ePC-SAFT--eNRTL comparison therefore measures the combined effect of the thermodynamic models and their reaction descriptions.
The two sets of equilibrium constants correspond to their respective reaction stoichiometries and activity references.

% CHECKLIST-BEGIN: method-02
% P:model-framework-11
The activity-coefficient description is defined by its extensive excess Gibbs energy $G^E(T,P,\mathbf n)$, with the eNRTL molecular and electrolyte interaction parameters specifying its composition dependence \cite{Zhang2011}.
On a mole-fraction reference basis,
\begin{align}
RT\ln\gamma_i^x&=\left(\pdv{G^E}{n_i}\right)_{T,P,n_{j\ne i}},\\
\mu_i&=\mu_i^{\circ,x}+RT\ln(x_i\gamma_i^x).
\label{eq:activity-chemical-potential}
\end{align}
For aqueous solutes, the conversion to the reaction system's molality reference preserves this chemical potential:
\begin{align}
\ln\!\left(\gamma_i^m\frac{m_i}{m^\circ}\right)
&=\ln(x_i\gamma_i^x)
+\frac{\mu_i^{\circ,x}-\mu_i^{\circ,m}}{RT}.
\label{eq:activity-reference-conversion}
\end{align}
Here $m_i$ is moles of solute per kilogram of water; water activity retains the solvent reference.
The same reference conversion is applied to the reaction constants, so changing the activity convention leaves the equilibrium condition $\boldsymbol\nu^{\mathsf T}\boldsymbol\mu=\mathbf0$ unchanged.
For molecular CO$_2$, the unsymmetric activity coefficient $\gamma_C^*=\gamma_C^x/\gamma_C^{x,\infty}$ gives $f_C^l=H_C^x x_C\gamma_C^*$.
The interaction records and model-specific reference offsets are specified in \cref{tab:comparison-thermo-inputs}.
The electrolyte-specific excess-Gibbs expression, long-range terms, and molecular/ion-pair interactions define the eNRTL comparator; the chemical-potential derivative identity alone does not specify those choices.
All paired liquid-model calculations use the same gas-side convention $f_j^v=y_j\phi_j^vP$; gas nonideality is therefore not changed when the liquid thermodynamic description is changed.
% CHECKLIST-END: method-02

\subsubsection{Fugacity and activity conventions}
% P:model-framework-12
At the gas--liquid interface, equality of chemical potential for a transferring molecular species is represented by
\begin{align}
f_{j,i}^v=f_{j,i}^l,\qquad j\in\{C,W\}.
\label{eq:interface-equilibrium}
\end{align}
For the gas mixture, $f_j^v=y_j\phi_j^vP$ and $\sum_jy_j=1$ over all four gas species.
For an activity-coefficient representation of the liquid,
\begin{align}
f_j^l=x_j\gamma_j f_j^\circ,
\end{align}
with the solute and solvent reference conventions explicitly distinguished.
The infinite-dilution Henry coefficient on a mole-fraction basis is
\begin{align}
H_j^x=\lim_{x_j\to0}\frac{f_j^l}{x_j}=\gamma_j^\infty f_j^\circ.
\end{align}
On a concentration basis, $f_j^l=H_j^c c_j$, so $H_j^x=\rho_l H_j^c$ in the corresponding dilute reference limit.
The distinction between $H^x$ and $H^c$ avoids introducing a density change merely by changing concentration units.

% P:model-framework-13
The transformation in \cref{eq:activity-reference-conversion} relates the eNRTL coefficients to the molality-based reaction equilibrium constants.
In ePC-SAFT, the corresponding quantities follow from the residual Helmholtz energy and the pressure-specified phase state.
Both descriptions therefore supply the same physical quantities to the absorber even though their free-energy models differ.

\subsubsection{Molecular and ionic contributions}
% P:model-framework-14
The hard-chain contribution connects the mixture hard-sphere reference to the chain lengths $m_i$,
\begin{align}
a^{\mathrm{hc}}=\overline m a^{\mathrm{hs}}
-\sum_i x_i(m_i-1)\ln g_{ii}^{\mathrm{hs}},\qquad
\overline m=\sum_i x_i m_i.
\end{align}
Here $g_{ii}^{\mathrm{hs}}$ is the hard-sphere contact radial distribution function.
Dispersion describes the attractive segment interactions through the segment diameter and energy parameters.
The association contribution accounts for the fractions $X_{iA}$ of unbonded sites,
\begin{align}
a^{\mathrm{assoc}}=\sum_i x_i\sum_A
\left(\ln X_{iA}-\frac{X_{iA}}2+\frac12\right).
\end{align}
Association-site topology and unlike-site interactions determine the molecular association pattern \cite{Gross2001,Cameretti2005}.
The electrolyte terms add ionic screening and solvation contributions; their dielectric and ion-size definitions belong to the selected thermodynamic parameterization.
For the neutral gas mixture, these ionic contributions vanish.
The pressure, chemical potentials, and composition derivatives must all be obtained from the same free-energy description.

\subsection{Enthalpy and thermal response}\label{subsec:calorics}
% P:model-framework-15
Each phase enthalpy is evaluated at its temperature, pressure, and composition with a consistent reference convention.
For a specified amount of reactive liquid, $H$ is the total enthalpy in joules and $C_{P,\mathrm{eq}}$ is in \unit{J\,K^{-1}}.
Define
\begin{align}
H_{\mathrm{eq}}(T,P,\boldsymbol\beta)
&=H\bigl(T,P,\mathbf n_{\mathrm{eq}}(T,P,\boldsymbol\beta)\bigr),\\
C_{P,\mathrm{eq}}
&=\left(\pdv{H_{\mathrm{eq}}}{T}\right)_{P,\boldsymbol\beta}.
\label{eq:equilibrium-heat-capacity}
\end{align}
The equilibrium derivative includes the change in species distribution with temperature.
It differs from the heat capacity at fixed composition.
Reaction and phase-change contributions enter through the species reference properties and mixture enthalpy; a separate constant heat of absorption would duplicate contributions already represented in this total enthalpy.
Phase enthalpy flows are evaluated on the same amount basis as the material balances.

\subsubsection{Energy references and temperature recovery}
% P:model-framework-16
An enthalpy reference fixes the additive energy assigned to each species at a reference temperature and pressure.
For a consistent reference heat capacity,
\begin{align}
h_i^\circ(T)=h_i^\circ(T_0)+\int_{T_0}^{T}c_{p,i}^\circ(\vartheta)\differential\vartheta.
\label{eq:reference-enthalpy}
\end{align}
Mixture enthalpy includes the thermodynamic contribution associated with the selected phase and composition.
The same reference convention is used for material crossing the interface, so interphase transfer conserves energy independently of the chosen zero of enthalpy.

% P:model-framework-17
For the reactive liquid, the temperature derivative can be resolved into a fixed-composition term and a speciation term,
\begin{align}
\left(\pdv{H_{\mathrm{eq}}}{T}\right)_{P,\boldsymbol\beta}
=\left(\pdv{H}{T}\right)_{P,\mathbf n}
+\left(\pdv{H}{\mathbf n}\right)_{T,P}
\left(\pdv{\mathbf n_{\mathrm{eq}}}{T}\right)_{P,\boldsymbol\beta}.
\end{align}
This distinction matters in the temperature bulge, where both the sensible energy and the equilibrium composition vary along the bed.

% P:model-framework-18
Writing the liquid enthalpy flow as $\dot H^l(T_l,P,\boldsymbol\beta_F)$, with $\boldsymbol\beta_F$ the mass/material-total flow vector and fixed zero charge flow, gives
\begin{align}
\odv{\dot H^l}{z}=
\left(\pdv{\dot H^l}{T_l}\right)_{P,\boldsymbol\beta_F}\odv{T_l}{z}
+\left(\pdv{\dot H^l}{P}\right)_{T_l,\boldsymbol\beta_F}\odv{P}{z}
+\left(\pdv{\dot H^l}{\boldsymbol\beta_F}\right)_{T_l,P}
\odv{\boldsymbol\beta_F}{z}.
\label{eq:enthalpy-total-derivative}
\end{align}
\Cref{eq:enthalpy-total-derivative} retains the composition contribution when recovering temperature from an energy balance.
For the constant-pressure bed, the middle term vanishes.
This relation also gives a common basis for comparing an enthalpy-state and temperature-state representation of the same physical column.

\subsection{Hydraulics and reference transfer model}\label{subsec:transport}
% P:model-framework-19
Packing geometry, liquid wetting, holdup, viscosity, diffusivity, and phase velocity determine the effective interfacial area $a_e$ and phase transfer coefficients.
The packed-column correlations follow the descriptions of Billet and Schultes, Tsai, and Chinen and co-workers \cite{Billet1999,Tsai2010,Chinen2018}.
For phase $p$, superficial velocity is $u_p=\dot V_p/A$, where $\dot V_p$ is the volumetric flow obtained from the phase amounts and density.
These correlations provide a common hydraulic and transport basis for each paired comparison.

% P:model-framework-transport-uncertainty
Uncertainty in viscosity and diffusivity propagates through the transport correlations to the predicted absorber response.
Liquid viscosity and diffusivity enter Reynolds and Schmidt numbers, film coefficients, and the Hatta number of the enhancement-factor model.
They can therefore change both integrated CO$_2$ capture and the position and magnitude of the liquid-temperature rise.
Interfacial area, holdup, gas-side diffusion, and interphase heat transfer provide additional sources of uncertainty \cite{Razi2013,Morgan2018,Morgan2020}.

% P:model-framework-transport-comparison
Using the same transport correlations in a thermodynamic comparison controls their specification, but does not remove their uncertainty or its interaction with equilibrium and heat release.
In the equilibrium-manifold film, species diffusivities set the mobility and the effective film thickness sets the transport distance.
Varying these inputs tests the dependence of the predicted response on their assigned values.
Such variations represent a sensitivity study; a probability distribution for prediction uncertainty additionally requires measured input uncertainties and their dependence.


\subsubsection{Interfacial area and phase holdup}
% P:model-framework-20
The effective area converts interfacial flux per unit area into transfer per unit bed volume.
The wetting correlation is
\begin{align}
a_e=1.42a_p\left[\frac{\rho_{l,m}}{\sigma_l}g^{1/3}
\left(\frac{u_l\varepsilon}{a_p}\right)^{4/3}\right]^{0.12},
\label{eq:effective-area}
\end{align}
where $a_p$ is specific packing area, $\rho_{l,m}$ is liquid mass density, and $\sigma_l$ is surface tension.
This is the selected packed-column wetting form based on the Tsai correlation \cite{Tsai2010}; the packing length per cross section is $a_p/\varepsilon$.
The liquid holdup $h_l$ and gas void fraction $h_v$ satisfy
\begin{align}
h_v=\varepsilon-h_l,
\end{align}
with packing void fraction $\varepsilon$.
Holdup relates the phase superficial velocities to the space available for flow and enters the packed-column transfer correlations.

\subsubsection{Mass- and heat-transfer coefficients}
% P:model-framework-21
The liquid-film correlation can be expressed as
\begin{align}
k_l=C_L\left(\frac{g\rho_{l,m}}{\mu_l}\right)^{1/6}
 d_h^{-1/2}\left(\frac{u_l}{a_p}\right)^{1/3}D_C^{1/2},
\label{eq:liquid-transfer-coefficient}
\end{align}
where $d_h$ is hydraulic diameter and $C_L$ is the packing coefficient \cite{Billet1999,Chinen2018}.
The gas-film coefficient is expressed on a fugacity basis so that multiplication by a fugacity difference produces a molar flux.
Reynolds, Schmidt, and Sherwood numbers make the property dependence explicit:
\begin{align}
\mathrm{Re}_p=\frac{\rho_{p,m}u_pd_h}{\mu_p},\qquad
\mathrm{Sc}_p=\frac{\mu_p}{\rho_{p,m}D_p},\qquad
\mathrm{Sh}_p=\frac{k_{p,c}d_h}{D_p}.
\end{align}
Here $k_{p,c}$ denotes a concentration-based transfer coefficient.
A pressure-based gas coefficient and a concentration-based gas coefficient are related through the gas-state concentration/fugacity convention; they are not interchangeable without that conversion.

% P:model-framework-22
Sensible heat exchange is represented by $U_T(T_v-T_l)$ per unit interfacial area.
The heat-transfer coefficient follows the heat/mass-transfer analogy used with the packed-column correlations.
The total interfacial energy flux is
\begin{align}
q=U_T(T_v-T_l)+\sum_{j\in\{C,W\}}J_j\overline h_{j,i}^{\mathrm{tr}},
\label{eq:interfacial-energy}
\end{align}
Here $q$ is positive for energy transfer from gas to liquid and has units of \unit{W\,m^{-2}}; $J_j$ is positive in the same direction.
The quantity $\overline h_{j,i}^{\mathrm{tr}}$ is a partial molar enthalpy in \unit{J\,mol^{-1}} on the selected interfacial phase convention, rather than an independently added heat of absorption.
The sensible term and transferred-material enthalpy must refer to the same side of the interface; changing that convention redistributes these two contributions while preserving $q$.
% CHECKLIST-BEGIN: method-03
% P:model-framework-23
The transferred-material term uses the gas-side enthalpy convention evaluated at the local bulk gas temperature:
\begin{align}
\overline h_{j,i}^{\mathrm{tr}}
&=\overline h_j^v(T_v,P,\mathbf y)
=\left(\pdv{(n_v h^v)}{n_j^v}\right)_{T_v,P,n_{k\ne j}^v}.
\label{eq:transferred-enthalpy}
\end{align}
For the common ideal-gas caloric description, this becomes
\begin{align}
\overline h_j^v(T_v)
&=\int_{T_\mathrm{ref}}^{T_v}c_{p,j}^v(T)\differential T,
\qquad T_\mathrm{ref}=\qty{298.15}{K},
\end{align}
with the gas reference enthalpies set to zero at $T_\mathrm{ref}$ and the corresponding liquid reference offsets retained in the liquid enthalpy.
The sensible contribution remains $U_T(T_v-T_l)$ with this convention.
The same total $q$ enters the two phase balances with opposite transfer directions; absorption heat is represented through the liquid enthalpy response and is not added again as a separate source.
% CHECKLIST-END: method-03
Water evaporation or condensation therefore affects both component flows and the thermal response.

\subsubsection{Enhancement-factor reference}

% P:model-framework-24
The enhancement-factor reference relates physical and reactive liquid-film resistance through $E$.
In concentration form its local CO$_2$ flux satisfies
\begin{align}
J_C=E k_l(c_{C,i}-c_{C,b})
=k_g(f_{C,b}^v-f_{C,i}),
\label{eq:enhancement-reference}
\end{align}
where $k_l$ has units of \unit{m\,s^{-1}} and $k_g$ has units of \unit{mol\,m^{-2}\,s^{-1}\,Pa^{-1}}.
The subscripts $i$ and $b$ denote interface and bulk values.
The enhancement relation follows the selected reaction-rate approximation \cite{Gaspar2015}; the interface concentration and fugacity are related by the thermodynamic model used in the comparison.
Water transfer and interphase heat transfer use the same respective closures within each enhancement/film pair.

% P:model-framework-25
The Hatta number compares the reaction and diffusion scales within the reference liquid film,
\begin{align}
\mathrm{Ha}=\frac{\sqrt{k_2c_{A,b}D_C}}{k_l},
\end{align}
where $k_2$ is the selected concentration-based pseudo-second-order rate coefficient.
For an MEA/water catalytic rate representation, $k_2=k_Ac_{A,b}+k_Wc_{W,b}$, with coefficient units consistent with the concentration basis.
The enhancement relation depends on both this reaction scale and the finite supply of reactive species \cite{Gaspar2015}.
% CHECKLIST-BEGIN: method-04
% P:model-framework-26
The reference uses the explicit finite-reactant expression
\begin{align}
R_+&=\frac{D_Ac_A}{2D_{A^+}c_{A^+}},&
R_-&=\frac{D_Ac_A}{2D_{AC^-}c_{AC^-}},&
\widehat E&=\frac{D_Ac_A}{2D_C\widetilde c_C},\\
E&=1+\frac{\mathrm{Ha}-1}{1+\mathrm{Ha}(R_++R_-+2)/\widehat E},
\label{eq:explicit-enhancement}
\end{align}
where $A^+$ and $AC^-$ denote protonated MEA and carbamate.
The concentration used in this reference expression is $\widetilde c_C=c_C/1.04542981654115$; the same fixed factor is retained in every paired reference calculation.
Both ionic diffusivity terms use the shared common-property ion correlation $D_{\mathrm{ion}}$ in the appendix; the species-resolved mobility values apply to the equilibrium-manifold film calculation.
For concentrations in \unit{mol\,m^{-3}}, the catalytic coefficients are
\begin{align}
k_A&=2.003\times10^4\exp(-4742.0/T),\\
k_W&=4.147\exp(-3110/T),
\end{align}
in \unit{m^6\,mol^{-2}\,s^{-1}}, giving $k_2$ in \unit{m^3\,mol^{-1}\,s^{-1}}.
Temperature is evaluated in kelvin.
The numerical evaluation restricts the resulting enhancement factor to $1\leq E\leq10^4$; the kinetic coefficients are not separately clipped.
% CHECKLIST-REVIEW: The active coefficients are the Luo-labeled pair in Enhancement_Factor.py, which overwrites the earlier Putta pair. The fixed divisor has no retained derivation/source locator. Preserve the implemented definition for comparison; confirm its rationale, kinetic source/domain, and any replacement with the absorber owner. See SOURCE_MAP.md.
% CHECKLIST-END: method-04

% P:model-framework-27
For a locally linear Henry relation, elimination of the interfacial concentration from \cref{eq:enhancement-reference} gives
\begin{align}
J_C=K_C(f_{C,b}^v-f_{C,b}^l),\qquad
\frac1{K_C}=\frac1{k_g}+\frac{H_C^c}{E k_l}.
\label{eq:series-resistance}
\end{align}
This series-resistance expression separates gas-film resistance from the chemically enhanced liquid-film resistance.
For a nonlinear thermodynamic relation, the interface equations retain the concentration-to-fugacity mapping explicitly.
Thus the thermodynamic comparison can change the equilibrium relation while preserving the reference kinetic and hydraulic assumptions.

\subsection{Equilibrium-manifold film transport}\label{subsec:reactive-film}
% P:model-framework-28
The local CO$_2$ film follows homogeneous reactive-equilibrium states at fixed $T$ and $P$.
For positive bulk and local loadings, the path coordinate is $\lambda=\ln(\alpha/\alpha_b)$, with bulk state $\lambda=0$ and interface state $\lambda_i$.
Species diffusivities $D_i$, true mole fractions $x_i$, and molar density $\rho_l$ define
\begin{align}
D_{ij}&=\frac{2D_iD_j}{D_i+D_j},\qquad w_{ij}=\rho_l x_i x_jD_{ij},\\
\mathbf L&=\operatorname{diag}(\mathbf w\mathbf1)-\mathbf w,\\
\mathbf L^P&=\mathbf L-\mathbf L\mathbf C^{\mathsf T}
(\mathbf C\mathbf L\mathbf C^{\mathsf T})^+\mathbf C\mathbf L.
\label{eq:film-projection}
\end{align}
The Laplacian mobility imposes zero total molar diffusion flux.
The rows of $\mathbf C$ impose zero electrical current and zero MEA- and water-component flux in the CO$_2$ film calculation; the superscript $+$ denotes the pseudoinverse.
With independent constraints on the zero-total-flux subspace, the projection is evaluated through the corresponding Gram-system linear solve; common scaling of both sides preserves the projection.
The constrained mobility is symmetric and positive semidefinite and has units of \unit{mol\,m^{-1}\,s^{-1}}.
The pair diffusivities are a mobility approximation whose physical meaning remains distinct from the thermodynamic driving force.

% P:model-framework-29
The chemical-potential tangent is
\begin{align}
\boldsymbol\tau(\lambda)
&=\odv{(\boldsymbol\mu/RT)}{\lambda}
=\left(\pdv{(\boldsymbol\mu_{\mathrm{eq}}/RT)}{\boldsymbol\beta}\right)_{T,P}
\odv{\boldsymbol\beta}{\lambda}.
\label{eq:film-tangent}
\end{align}
The loading-to-total mapping and the equilibrium derivative use the same complete conserved basis.
For the normalized apparent feed $\widehat{\mathbf n}_{\mathrm{feed}}$, changing the CO$_2$ amount by $\exp(\lambda)$ while holding the MEA and water amounts fixed gives
\begin{align}
\odv{\widehat{\mathbf n}_{\mathrm{feed}}}{\lambda}
&=\widehat n_{C,\mathrm{feed}}\left(\mathbf e_C-\widehat{\mathbf n}_{\mathrm{feed}}\right),\\
\odv{\boldsymbol\beta}{\lambda}
&=\begin{pmatrix}\mathbf M^{\mathsf T}\\\mathbf B\end{pmatrix}
\odv{\widehat{\mathbf n}_{\mathrm{feed}}}{\lambda}.
\label{eq:feed-loading-direction}
\end{align}
The fixed-temperature, fixed-pressure and fixed-zero-charge directions vanish.
This is a feed-to-invariant derivative; the derivative of the true equilibrium species amounts is supplied by the equilibrium response, not by identifying those amounts with the normalized feed.
For the CO$_2$ component vector $\mathbf b_C$, the liquid-film integral and gas-film balance give
\begin{align}
J_C&=\frac{1}{\delta}\int_0^{\lambda_i}
\mathbf b_C^{\mathsf T}\mathbf L^P(\lambda)\boldsymbol\tau(\lambda)\differential\lambda
=k_g\left[f_{C,b}^v-f_{C,i}(\lambda_i)\right],
\label{eq:film-interface}\\
\delta&=\frac{D_C^{\mathrm{column}}}{k_l}.
\label{eq:film-thickness}
\end{align}
Positive $J_C$ denotes transfer into the liquid.
Quadrature evaluates the equilibrium-path integral, and a scalar root determines the interfacial loading.
The path-integrated resistance replaces the enhancement multiplier in \cref{eq:enhancement-reference}.
The column diffusivity used to define $\delta$ and the species diffusivities in $\mathbf L$ retain their separate roles.

\subsection{Column balances and boundary conditions}\label{subsec:governing-equations}
% P:model-framework-30
Let $J_j$ denote interfacial component flux into the liquid and $q$ the total interfacial energy flux in the same direction.
For an adiabatic bed, the countercurrent balances are
\begin{align}
\odv{F_j^l}{z}&=-Aa_eJ_j,&
\odv{F_j^v}{z}&=-Aa_eJ_j,\qquad j\in\{C,W\},
\label{eq:material-balances}\\
\odv{\dot H^l}{z}&=-Aa_eq,&
\odv{\dot H^v}{z}&=-Aa_eq,\qquad \odv{P}{z}=0.
\label{eq:energy-balances}
\end{align}
Here $W$ denotes the water component; $q$ includes sensible heat exchange and the enthalpy carried by transferred material, evaluated with a common interfacial convention.
The signs in \crefrange{eq:material-balances}{eq:energy-balances} reflect the opposing phase-flow directions.
Consequently, $F_j^v-F_j^l$ and $\dot H^v-\dot H^l$ are constant along the bed.

% P:model-framework-31
One energy-coordinate representation of the physical balances is
\begin{align}
\mathbf y=(F_C^l,F_W^l,F_C^v,F_W^v,\dot H^l,\dot H^v,P)^{\mathsf T}.
\end{align}
Gas inlet component flows and enthalpy are specified at $z=0$; liquid inlet component flows and enthalpy are specified at $z=H$.
The remaining nontransferring component flows are fixed by the feeds, and the inlet pressure closes the pressure state.
At each axial position, the phase amounts, equilibrium composition, and enthalpy determine the local temperatures and transfer rates.
The numerical unknowns need not be these enthalpy coordinates: temperatures, scaled component flows, and interfacial variables can parameterize the same physical balance quantities.
The enthalpy chain rule in \cref{eq:enthalpy-total-derivative} or the conservative balance map in \cref{eq:conservative-trapezoid} retains the energy equation under such a change of coordinates.

\subsubsection{Integral conservation and column observables}
% P:model-framework-32
Integrating the CO$_2$ balance over the bed gives
\begin{align}
F_C^v(0)-F_C^v(H)
=F_C^l(0)-F_C^l(H)
=A\int_0^H a_e(z)J_C(z)\differential z.
\label{eq:integral-uptake}
\end{align}
The same relation applies to water with its signed interfacial flux.
For the adiabatic column,
\begin{align}
\dot H^v(0)+\dot H^l(H)=\dot H^v(H)+\dot H^l(0).
\label{eq:integral-energy}
\end{align}
These identities connect the local flux and energy laws to inlet/outlet measurements.
They also make the countercurrent sign convention explicit: absorption increases liquid CO$_2$ flow toward the bottom while decreasing gas CO$_2$ flow toward the top.

% P:model-framework-33
The cumulative uptake to height $z$ is
\begin{align}
\mathcal U_C(z)=A\int_0^z a_e(\zeta)J_C(\zeta)\differential\zeta
=F_C^v(0)-F_C^v(z).
\end{align}
A normalized profile $\mathcal U_C(z)/\mathcal U_C(H)$ distinguishes the distribution of absorption from its total magnitude.
The temperature maximum and its axial position provide a corresponding description of the thermal response.
These local and integral measures are used together in the model comparisons.
